% 默认入口：本文件等同于 main-cumcmthesis.tex，使用 cumcmthesis 类。
% 如果本机没有 cumcmthesis.cls，请改用 main-ctexart-fallback.tex。

\documentclass[withoutpreface,bwprint]{cumcmthesis}

\usepackage{amsmath,amssymb,booktabs,graphicx}
\usepackage{hyperref}
\usepackage{url}

% Reference mode: set \userefbibtrue to use ref.bib with BibTeX.
\newif\ifuserefbib
\userefbibfalse

\title{论文题目占位}
\tihao{A}
\baominghao{000000000000}

\begin{document}

\maketitle

\begin{abstract}
这里撰写摘要。摘要应概括问题、模型、算法、结果、检验与结论。不得出现可识别参赛者或机构的真实身份信息。

\keywords{关键词一\quad 关键词二\quad 关键词三}
\end{abstract}

% CUMCMThesis 示例说明 2019 年起电子版通常不需要目录。
% 若当年官方规则明确要求目录，再启用下面两行。
% \tableofcontents
% \newpage

\section{问题重述}
说明问题背景、研究对象和需要完成的任务。

\section{问题分析}
分析问题的关键因素、数据特征和建模思路。

\section{模型假设}
\begin{itemize}
  \item 假设一：说明假设内容及合理性。
  \item 假设二：说明假设内容及合理性。
\end{itemize}

\section{符号说明}
\begin{table}[htbp]
  \centering
  \caption{主要符号说明}
  \begin{tabular}{cl}
    \toprule
    符号 & 含义 \\
    \midrule
    $x$ & 决策变量占位 \\
    $t$ & 时间指标占位 \\
    \bottomrule
  \end{tabular}
\end{table}

\section{模型建立与求解}
\subsection{模型一}
给出模型结构、参数估计、算法步骤和求解结果。

\begin{equation}
  \min_x f(x), \quad \text{s.t.}\quad g_i(x)\le 0,\ i=1,\ldots,m.
\end{equation}

\subsection{模型二}
根据题目需要扩展、修正或对比模型。

\section{模型检验与灵敏度分析}
说明模型验证方法，并分析关键参数变化对结果的影响。

\section{模型评价与推广}
总结模型优点、局限性和可推广方向。

\section{结论}
凝练给出最终结论和建议。

\ifuserefbib
\nocite{placeholder-ref}
\bibliographystyle{unsrt}
\bibliography{ref}
\else
\begin{thebibliography}{99}
\bibitem{placeholder} 作者. 文献题名占位[J]. 期刊或出版社, 年份.
\end{thebibliography}
\fi

\appendix
\section{附录：补充材料}
放置推导、算法伪代码、扩展表格或图形说明。

\section{支撑材料文件列表}
\begin{itemize}
  \item \texttt{code/main.py}：主要求解程序占位。
  \item \texttt{data/input.csv}：输入数据占位。
  \item \texttt{results/output.csv}：输出结果占位。
\end{itemize}

\section{AI 工具使用详情说明}
如使用 AI 工具，应按照当年竞赛规则如实说明工具名称、版本或访问日期、使用环节、生成内容范围和人工核验方式。如未使用，应按当年规则填写未使用说明。

\end{document}
